%% ma: L10-B02-0033
%% lop: 10
%% chuong: 1
%% bai: 2
%% dang: 10.2.4
%% loai: TN4
%% muc: TH
%% dap_an: "B"
%% nguon: "Ảnh giáo viên gửi 10/10/2026 (ThS. Nguyễn Hoàng Việt, Chương 2 Tập hợp), Đề 2 (đề thứ hai, không ghi tiêu đề), Phần 1 Câu 10"
%% kien_thuc: "Bài 2 (tập hợp và các phép toán trên tập hợp)"
%% trang_thai: cho-duyet
%% ly_do: "Dạng toán mới đề xuất 10.2.4: tập con và số tập con"
\begin{ex}
Cho tập $A=\{x\in\mathbb N\mid(2-x)(x^2-3x-4)=0\}$. Hỏi tập $A$ có bao nhiêu tập con?
\choice{$2$}{\True $4$}{$7$}{$8$}
\loigiai{$2-x=0\Leftrightarrow x=2$; $x^2-3x-4=(x-4)(x+1)=0\Leftrightarrow x\in\{4;-1\}$. Nghiệm tự nhiên: $A=\{2;4\}$ có $2$ phần tử nên có $2^2=4$ tập con.}
\end{ex}
